\chapter{Compilazione ed esecuzione del sistema}
\section{Requisiti}
I requisiti per eseguire il sistema sono i seguenti:
\begin{itemize}
    \item Erlang/OTP dalla versione 25;
    \item Python dalla versione 3.5 in poi;
    \item le librerie python Flask e PyVis.
\end{itemize}
Per installare le librerie python basta eseguire i comandi
\begin{verbatim}
    pip install flask
    pip install pyvis
\end{verbatim}

\section{Compilazione}
La compilazione per l'interfaccia grafica non è necessaria, in quanto Python è un linguaggio interpretato.
Per il resto del sistema, è possibile compilare tutto l'\texttt{Emakefile} offerto, usando il comando \texttt{erl -make}.

\section{Esecuzione}
Per eseguire l'interfaccia grafica basta usare il comando \texttt{flask -app app run}

Per eseguire il sistema ACO è necessario utilizzare diverse chiamate.
La prima da eseguire è \texttt{user\_api:system\_setup()} o \texttt{user\_api:system\_setup\_simple()}, se si vuole vuole avere
il processo \texttt{user\_node} o meno. 
Appena tutti i nodi distribuiti sono pronti, è necessario connetterli (ad esempio tramite \texttt{net\_adm:ping()}) e poi avviare
l'elezione del ticket master con \texttt{user\_api:start\_election()}.

Ora è necessario aggiungere gli uffici del sistema, per farlo bisogna usare \texttt{user\_api:add\_office()}, che richiede un 
nome per l'ufficio e l'id del nodo distribuiti a cui afferisce.

Dopodiché, si possono usare \texttt{user\_api:init\_add\_clients} e \texttt{user\_api:init\_add\_costs} per aggiungere nodi e archi.
Il primo metodo richiede una lista di coppie, formate da una lista di nodi e il nome di un ufficio. Il secondo metodo richiede una lista 
di tripla, dove i primi due elementi sono i nodi connessi dall'arco e il terzo è il costo dell'arco.
Ultimo passaggio è quello di avviare l'algoritmo ACO, usando la chiamata \texttt{user\_api:system\_start}
Non è possibile avere più processi \texttt{flooding} sulla stessa macchina, in quanto andrebbero in conflitto scrivendo e leggendo
sullo stesso file.

Dopo aver 3 terminali aperti, un possibile esempio dei passaggi da eseguire:
\begin{verbatim}
    # ad ogni step guardare la pagina web se si è aggiornata

    # sul terminale 1
    cd WebInterface
    flask --app app run

    # sul terminale 2
    erl -make

    # sul  terminale 2
    erl -sname Utente
    # sul terminale  3
    erl -sname Utente2

    # su un nodo che deve contenere il nodo utente:
    # possiamo scegliere il terminale 2 o 3
    user_api:system_setup().

    # sugli altri nodi: 
    user_api:system_setup_simple().

    # su tutti i nodi che non sono l'utente2
    net_adm:ping('Utente2@Computer').

    # aggiungere tutti gli uffici nel nodo aco che contiene il nodo utente
    user_api:add_office(milano, node()).
    user_api:add_office(udine, 'Utente2@Computer').
    user_api:add_office(padova, 'Utente2@Computer').
    # in questo caso il nostro nodo contiene solo l'ufficio di Milano e l'altro nodo
    # contiene quelli di Udine e padova

    # facciamo partire l'elezione
    user_api:start_election().

    # sul nodo aco che contiene il nodo utente
    # aggiungiamo 4 nodi, 2 a milano, uno a udine e uno a padova
    user_api:init_add_clients([{[1,2], milano}, {[3], udine}, {[4], padova}]).

    # aggiungiamo i costi relativi ai nodi
    user_api:init_add_costs([{1, 2, 5}, {1, 3, 6}, {1, 4,7}, 
            {2, 3, 8}, {2, 4, 9}, {3, 4, 10}]).

    # eseguiamo l'algoritmo 
    user_api:system_start().

    # aggiungiamo due clienti in corso durante l'algoritmo
    user_api:add_client(5, udine, [{1, 5, 1}, {2, 5, 1},
                 {3, 5, 1}, {4, 5, 1}]).

    user_api:add_client(6, padova, [{1, 6, 1}, {2, 6, 1}, {3, 6, 1}, 
            {4,6, 1}, {5,6, 1}]).

    # rimozione dei due clienti in corso
    user_api:remove_client(6).

    user_api:remove_client(5).

    # aggiorniamo i costi degli archi
    user_api:update_costs([{1,2,1}, {1,3,1}, {1,4,1}, 
            {2,3, 1}, {2,4,1}, {3,4, 1}]).

    # ora chiudiamo il terminale 3
    # riapriamo il terminale 3
    # nel terminale 2 riconnettiamo il terminale 3
    net_adm:ping('Utente2@Computer').

    # eseguiamo nel terminale 3:
    user_api:system_setup_simple().
    ticket_system:connect_to_system().
    rr(aco_node).

    # aggiungiamo il nodo 3 localmente sul terminale 3
    aco_node:add_local_node(#node{name=3}).
     
    # aggiungiamo gli archi locali al nodo 3
    aco_node:add_edges([#edge{from=#node{name=1}, to=#node{name=3}, weight=6}, 
    #edge{from=#node{name=2}, to=#node{name=3}, weight=8}]). 
    
    # aggiungiamo il nodo aco al sistema, questa chiamata potrebbe impiegarci un attimo
    aco_node:running_init(3).

    # tornando sul terminale 2 eseguiamo questo
    user_api:add_client(4, udine, [{1, 4, 7}, {2, 4, 9},
                 {3, 4, 10}]).

    # il sistema quindi supporta l'aggiunta, la rimozione di aco e cliente
    # e l'aggiornamento di archi 

    # il grafo deve aggiornasi dopo aver perso il secondo nodo aco
\end{verbatim}
